\section{Beschreibung des Versuchsdesigns}
\label{ch:Ergebnisse:sec:Versuchsdesign}

%Forschungsfrage oder Hypothese: Was soll untersucht werden? 
Es wird untersucht, inwiefern sich die Ausprägung der Usability"=Kriterien der Effektivität, Effizienz und Zufriedenheit verändert, wenn für den Usability"=Test zum einen das alte und zum anderen das neue System verwendet werden und wie der Nutzer nach diesem Test das jeweilige System bewertet.

%Unabhängige Variable(n) (Faktoren): Die Variablen, die vom Forscher manipuliert werden (z.B. Temperatur, Konzentration).
Die Daten teilen sich in abhängige und unabhängige Variablen. Die unabhängigen Variablen sind die Ursache für die zu erhebenen Daten und umfassen die einzelnen Fragen und Aufgaben. Jede Frage bzw. Aufgabe ist jeweils im alten und im neuen System von den Probanden zu bearbeiten.

%Abhängige Variable(n): Die Variablen, deren Werte von den unabhängigen Variablen abhängen (z.B. Reaktionszeit, Ausbeute). 
Die erhobenen Ergebnisse sind die drei abhängigen Variablen. Diese umfassen das subjektive Maß der Zufriedenheit sowie die objektiven Maße der Zeit und des Fehlers.
Als Zufriedenheit werden verschiedene Aspekte der Systeme bewertet. Die Zeit meint die Dauer, welche die Testperson zur Erfüllung der Aufgaben benötigt. Die Fehler werden als Rate oder Fehlerquote gemessen, zu wie viel Prozent die Aufgaben richtig beantwortet wurden.
Im Anschluss können die Zeit und die Fehlerquote zu einer Effizienzbewertung verrechnet werden.

%Störvariable(n): Variablen, die den Zusammenhang zwischen unabhängigen und abhängigen Variablen beeinflussen können, aber nicht das Hauptinteresse des Experiments sind (z.B. Raumtemperatur, Alter der Versuchsteilnehmer). 
Damit die Beeinflussung der Testpersonen gering gehalten wird und ein \textit{Reihenfolgeeffekt} \cite{AnnaSchantl:2009:Reihenfolgeeffekt} ausgeschlossen werden kann, werden die Probanden in zwei Versuchsgruppen aufgeteilt, welche die Systeme in unterschiedlicher Reihenfolge bearbeiten.

%Versuchsgruppe(n) oder Bedingungen: Die verschiedenen Kombinationen von Faktorstufen, die im Experiment untersucht werden. 


%Versuchsplan (Design): Ein strukturierter Plan, der festlegt, wie viele Versuchspunkte durchgeführt werden und wie die Versuchsgruppen oder Bedingungen randomisiert werden. 


%Datenerhebung: Wie und wann werden Daten gesammelt. 
Jeder Versuch wird aufgezeichnet, daher können die Daten kontrolliert und wiederholbar erfasst werden.

%Datenanalyse: Welche statistischen Methoden werden verwendet, um die Daten auszuwerten. 
Um die Daten auszuwerten, werden die Mittelwerte zur Bestimmung der Tendenz des Ergebnisses verglichen sowie der t-Test (s. \autoref{ch:Ergebnisse:sec:Hypothesen_Auswertung:sub:Bsp_TTest}) zur Bestimmung der statistischen Signifikanz der Aussage durchgeführt.

%Interpretation der Ergebnisse: Was bedeuten die Ergebnisse und welche Schlussfolgerungen können gezogen werden. 
Alle aufgezeichneten Daten sind metrisch und es kann eindeutig interpretiert werden, welches System für diesen Anwendungsfall besser geeignet sein wird.


%-	abhängige Variablen (Wirkung / Ergebnis):\\
%->	subjektives Maß: Usability\\
%->	objektives Maß: Zeit und Fehler\\
%- zu Effizienz verrechnen\\
%-> unabhängige Variablen (Ursache)\\
%-> das was ich manipuliere: alt-neu (2-fachgestuft)\\
%->	Aufgabenbeschreibung
